\فصل{بحث، نتیجه‌گیری و پیشنهادها}

در این فصل، ابتدا خلاصه‌ای از یافته‌های پژوهش ارائه می‌شود.
سپس یافته‌ها در مقایسه با ادبیات موضوع مورد بحث و تفسیر قرار می‌گیرند.
در ادامه، سهم‌های علمی پایان‌نامه مرور شده
و محدودیت‌های پژوهش بیان می‌شوند.
در نهایت، پیشنهادهایی برای پژوهش‌های آتی ارائه می‌شود.


\قسمت{خلاصه یافته‌ها}

در این پایان‌نامه، مسئله زمان‌بندی کارگاهی در بستر یک زیرساخت رایانشی سه‌لایه‌ای
لبه-مه-ابر با سه تابع هدفِ زمان تکمیل کل، مصرف انرژی و عدم‌قابلیت‌اطمینان
(معادلِ بیشینه‌سازی قابلیت اطمینان) مدل‌سازی و حل شد.
الگوریتم ژنتیک بهبودیافته مبتنی بر \lr{NSGA-II} با سه مؤلفه بهبودی
(مقداردهی اولیه هوشمند، انتخاب تطبیقی عملگر مبتنی بر \lr{UCB1}
و جست‌وجوی محلی) پیشنهاد و پیاده‌سازی شد.

یافته‌های اصلی این پژوهش به شرح زیر خلاصه می‌شوند:

\شروع{شمارش}
\فقره
\مهم{عملکرد برتر الگوریتم بهبودیافته}:
الگوریتم ژنتیک بهبودیافته در اکثر نمونه‌های معیار از نظر شاخص ابرحجم
و فاصله نسلی معکوس، عملکرد بهتری نسبت به الگوریتم ژنتیک ساده
و قواعد ابتکاری \lr{SPT} و \lr{MWR} نشان داد.
آزمون فریدمن تفاوت معنادار آماری میان الگوریتم‌ها را تأیید کرد
و آزمون‌های ویلکاکسون با تصحیح هولم نیز برتری الگوریتم بهبودیافته
را در اکثر مقایسه‌های دودویی نشان دادند.

\فقره
\مهم{نقش مؤلفه‌های بهبودی}:
مطالعه حذفی نشان داد که هر دو مؤلفه بررسی‌شده سهم مثبتی
در عملکرد نهایی الگوریتم دارند.
حذف جست‌وجوی محلی بیشترین تأثیر منفی را بر شاخص ابرحجم داشت.
حذف انتخاب تطبیقی عملگر نیز
افت عملکرد قابل‌ملاحظه‌ای ایجاد کرد.

\فقره
\مهم{ارزیابی در بستر زیرساخت ناهمگن}:
مدل‌سازی مسئله در بستر لبه-مه-ابر با شبیه‌سازی \lr{YAFS}
نشان داد که در نظر گرفتن ناهمگنی زیرساخت (تفاوت توان پردازشی، مصرف انرژی
و نرخ خرابی گره‌ها) تأثیر معناداری بر کیفیت راه‌حل‌های پارتو دارد.

\فقره
\مهم{قابلیت بازتولید}:
تمامی آزمایش‌ها با ۱۰ بذر تصادفی ثابت روی ۱۵ نمونه معیار اجرا شدند
و نتایج با فاصله اطمینان ۹۵٪ گزارش شدند.
کدها و پیکربندی‌های لازم به‌منظور بازتولید نتایج مستندسازی شده‌اند.
\پایان{شمارش}


\قسمت{بحث و تفسیر نتایج}

یافته‌های این پژوهش با ادبیات موجود همسو است و چند نکته کلیدی را روشن می‌سازد.

نخست، برتری الگوریتم ژنتیک بهبودیافته نسبت به نسخه ساده
تأییدکننده اثربخشی رویکرد ترکیب مؤلفه‌های بهبودی در الگوریتم‌های تکاملی است.
این نتیجه با یافته‌های پژوهش‌های پیشین
در زمینه الگوریتم‌های میمتیک\پاورقی{Memetic Algorithms}
و ترکیب جست‌وجوی محلی با الگوریتم‌های جمعیتی سازگار است.

دوم، مطالعه حذفی نشان داد که جست‌وجوی محلی بیشترین سهم
را در بهبود عملکرد دارد.
این امر نشان می‌دهد که بهره‌برداری\پاورقی{Exploitation}
از فضای اطراف راه‌حل‌های خوب
در مسائل زمان‌بندی ترکیبیاتی اهمیت ویژه‌ای دارد.
از سوی دیگر، سایر مؤلفه‌های طراحی الگوریتم
با تسریع همگرایی اولیه و تنظیم پویای تعادل اکتشاف و بهره‌برداری،
سهم مکمل مهمی ایفا می‌کنند.

سوم، مدل‌سازی مسئله در بستر زیرساخت ناهمگن لبه-مه-ابر نشان داد
که نگاشت ماشین‌ها به گره‌های رایانشی
تأثیر قابل‌توجهی بر موازنه\پاورقی{Trade-off} بین اهداف دارد.
گره‌های ابری با توان بالا زمان تکمیل کمتری ارائه می‌دهند
اما مصرف انرژی بیشتری دارند،
در حالی که گره‌های لبه‌ای مصرف انرژی کمتری دارند
اما نرخ خرابی بالاتری نشان می‌دهند.
این مبادله، ضرورت به‌کارگیری رویکرد چندهدفه
را به‌خوبی نمایان می‌سازد.


\قسمت{سهم‌های علمی پژوهش}

سهم‌های اصلی این پایان‌نامه به‌صورت خلاصه عبارتند از:

\شروع{شمارش}
\فقره
ارائه مدل سه‌هدفه برای مسئله زمان‌بندی کارگاهی
در بستر زیرساخت رایانشی لبه-مه-ابر
با مدل‌های صریح انرژی و قابلیت اطمینان.
\فقره
طراحی و پیاده‌سازی الگوریتم ژنتیک بهبودیافته
با ترکیب مقداردهی اولیه هوشمند (مبتنی بر قواعد \lr{SPT} و \lr{MWR})،
انتخاب تطبیقی عملگر (مبتنی بر \lr{UCB1})
و جست‌وجوی محلی در چارچوب \lr{NSGA-II}.
\فقره
ارزیابی تجربی جامع و قابل بازتولید
روی ۱۵ نمونه معیار استاندارد
با آزمون‌های آماری ناپارامتری و مطالعه حذفی.
\پایان{شمارش}


\قسمت{محدودیت‌های پژوهش}

این پژوهش دارای محدودیت‌های زیر است
که لازم است در تفسیر نتایج مورد توجه قرار گیرند:

\شروع{فقرات}
\فقره
زیرساخت شبیه‌سازی‌شده ایستا فرض شده
و تغییرات پویای شبکه (مانند خرابی گره‌ها حین اجرا
یا تغییر پهنای باند) در نظر گرفته نشده است.
\فقره
تنها مسئله زمان‌بندی کارگاهی کلاسیک بررسی شده
و حالت‌های عمومی‌تر مانند
زمان‌بندی کارگاهی منعطف\پاورقی{Flexible JSSP}
یا زمان‌های آماده‌سازی وابسته به توالی
مدل‌سازی نشده‌اند.
\فقره
اندازه جمعیت و تعداد نسل‌ها برای تمامی نمونه‌ها ثابت بوده
و تنظیم خودکار ابرپارامترها بررسی نشده است.
\فقره
مقایسه‌ها تنها با الگوریتم ژنتیک ساده و قواعد ابتکاری انجام شده
و سایر روش‌های مرجع
به‌عنوان پایه مقایسه استفاده نشده‌اند.
\فقره
نتایج مبتنی بر شبیه‌سازی بوده
و اعتبارسنجی در محیط واقعی تولیدی یا صنعتی انجام نشده است.
\پایان{فقرات}


\قسمت{پیشنهادها برای پژوهش‌های آتی}

بر اساس نتایج و محدودیت‌های این پژوهش،
پیشنهادهای زیر برای ادامه پژوهش ارائه می‌شوند:

\شروع{شمارش}
\فقره
\مهم{تعمیم به مسئله زمان‌بندی کارگاهی منعطف}:
تعمیم مدل پیشنهادی به مسئله زمان‌بندی کارگاهی منعطف
که در آن هر عملیات می‌تواند روی بیش از یک ماشین اجرا شود.
\فقره
\مهم{مدل‌سازی زیرساخت پویا}:
در نظر گرفتن تغییرات پویای شبکه
مانند خرابی و بازیابی گره‌ها، تغییر بار شبکه و مهاجرت وظایف
و طراحی سازوکار زمان‌بندی مجدد\پاورقی{Rescheduling}.
\فقره
\مهم{مقایسه با الگوریتم‌های بیشتر}:
مقایسه با سایر الگوریتم‌های چندهدفه
به‌منظور ارزیابی جامع‌تر.
\فقره
\مهم{بهره‌گیری از یادگیری تقویتی عمیق}:
استفاده از روش‌های یادگیری تقویتی عمیق\پاورقی{Deep Reinforcement Learning}
برای انتخاب تطبیقی عملگر به‌جای سازوکار \lr{UCB1}
با هدف بهبود قابلیت تعمیم‌پذیری.
\فقره
\مهم{تنظیم خودکار ابرپارامترها}:
بررسی روش‌های تنظیم خودکار\پاورقی{Auto-Tuning}
مانند بهینه‌سازی بیزی\پاورقی{Bayesian Optimization}
برای یافتن پیکربندی بهینه ابرپارامترهای الگوریتم ژنتیک
متناسب با ابعاد هر نمونه مسئله.
\فقره
\مهم{ارزیابی روی نمونه‌های بزرگ‌تر}:
آزمایش الگوریتم روی نمونه‌های مسئله با ابعاد بسیار بزرگ
(بیش از ۱۰۰ کار و ۲۰ ماشین)
و بررسی مقیاس‌پذیری\پاورقی{Scalability} رویکرد پیشنهادی.
\پایان{شمارش}


\قسمت{جمع‌بندی فصل}

این فصل، جمع‌بندی نهایی پژوهش را با تأکید بر دستاوردها، محدودیت‌ها و مسیرهای آینده ارائه داد.
بر پایه نتایج فصل پیشین، می‌توان نتیجه گرفت که ترکیب راهکارهای بهبودی پیشنهادی
در چارچوب \lr{NSGA-II} برای حل چندهدفه \lr{JSSP} در زیرساخت ناهمگن، رویکردی مؤثر و قابل توسعه است.
